\chapter{Library organization and Lean interfaces}

The file organization follows the mathematical dependency plan. The complete
node-to-file assignment and import graph are recorded in
\href{modules.json}{the module manifest}; the
\href{modules.svg}{module graph} records the proposed imports.
This is a proposed architecture, checked for cycles. A file appearing there
need not yet contain a compiled statement. Source validation and Lean build
results are reported separately.

\section{Principles}
Definitions carry only the structure needed to state the notion. Local
finiteness, positivity of dimension, measurability, and finite Hausdorff mass
are theorem hypotheses when required. The library reuses mathlib's measures,
Hausdorff measure, Euclidean spaces, continuous multilinear maps, projections,
and Gaussian distributions. It reuses FormalizingGMT's density and tangent
predicates and compact-test convergence; convention changes have named adapters.

This follows mathlib's guidance on
\href{https://mathlib.org/contribute/style.html}{file organization and style}
and \href{https://mathlib.org/contribute/naming.html}{declaration naming}:
put reusable results in their mathematical topic, use meaningful theorem names,
and keep imports narrow. The initial namespace is \texttt{Preiss}; moving a
result upstream should not require changing its mathematical hypotheses.

\section{File families}
\begin{itemize}
\item \texttt{Geometry}: Euclidean spaces, planes, projections, compact plane
parameters, cones, Lipschitz extension, and Hausdorff measure on planes.
These files do not import tangent-measure or uniform-measure theory.
\item \texttt{Rectifiability/Defs} and \texttt{GraphDefs}: partial-domain
Lipschitz pieces, exact countable coverage, rectifiability modulo null sets,
rectifiable measures, and Lipschitz graphs. Basic carrier calculus and
rectifiable/purely unrectifiable decomposition build on these definitions.
\item \texttt{Density}: conventions, power-law comparisons with Hausdorff
measure, finite carrier pieces, and the zero/excess-dimension endpoints.
\item \texttt{Tangent}: dimensional scaling, normalization adapters,
compactness, restrictions, recentering, and closure under further blow-ups.
Generic scaling algebra is separate from the final flat-pseudotangent theorem.
\item \texttt{Analysis} and \texttt{MeasureTheory}: graded projections,
polynomial feature maps, spectral optimization, vague convergence, and
Gaussian integrability/determination. These statements should be useful
without assuming that a measure is uniform.
\item \texttt{Uniform}: unit ball laws, radial integrals, moment forms,
Taylor coefficients, blow-downs, conical gaps, and flat-at-infinity rigidity.
Coefficient families are explicit data; their asymptotic properties are
hypotheses or conclusions of named theorems.
\item \texttt{Rectifiability/WeakApproximation} and
\texttt{Rectifiability/MarstrandMattila}: approximation, restriction,
projection estimates, packing, and the sufficient rectifiability criterion.
\item The companion graph-cover, graph-area, and tangent-necessity modules
prove the converse and full characterization. They are outside the main
Preiss theorem's import closure.
\item \texttt{Main} assembles the two tangent arguments, the sufficient
criterion, and the endpoint dimensions. It introduces no new geometric notion.
\end{itemize}

\section{Statement and proof status}
Each theorem is introduced in its canonical topic file with an explicit
\texttt{sorry} body, then proved in place. The earlier \texttt{PreissBlueprint} statement laboratory has been replaced by
a compatibility import of the canonical library.
Concrete definitions must not contain placeholders. Compiling all statements
is a separate milestone from proving them: an assembly theorem may have a
written proof but still depend transitively on \texttt{sorryAx}.

Pure local modules use Lean's modern module system. FormalizingGMT currently
uses legacy modules, so files importing it remain explicit legacy boundaries;
a blanket conversion would be rejected by the compiler. This is a build
constraint, not an additional mathematical dependency.
